%% hydroflux methods paper for Environmental Modelling & Software (EMS)
%% Using elsarticle with Harvard (author-year) bibliography style.
%%
%% Compile with xelatex (native unicode in monospace + math fonts):
%%   xelatex paper && bibtex paper && xelatex paper && xelatex paper
%%
\documentclass[preprint,review,12pt,authoryear]{elsarticle}

%% --- Font / unicode setup for xelatex ---
\usepackage{fontspec}
\setmainfont{Latin Modern Roman}
\setsansfont{Latin Modern Sans}
\setmonofont[Scale=0.95]{DejaVu Sans Mono}
\usepackage{unicode-math}
\setmathfont{Latin Modern Math}

%% --- Math, tables, references, code ---
\usepackage{amsmath}
\usepackage{booktabs}
\usepackage{calc}
\usepackage{longtable}
\usepackage{multirow}
\usepackage{array}
% pandoc helpers for longtable column widths
\newlength{\hfdcw}
\setlength{\hfdcw}{\columnwidth}
\providecommand{\real}[1]{#1}
\usepackage{lineno}
\usepackage{hyperref}
\usepackage{xcolor}
\usepackage{graphicx}
\usepackage{siunitx}

\journal{Environmental Modelling \& Software}

\begin{document}

\begin{frontmatter}

\title{hydroflux: A Well-Balanced, Differentiable-by-Design 2D
Shallow-Water Solver in Rust, Verified Against Analytical and Community
Benchmarks and Applied to a Semiarid Andean Reach}

\author[usach]{Francisco Parra\corref{cor1}}
\ead{francisco.parra.o@usach.cl}
\cortext[cor1]{Corresponding author}

\author[unsl]{Verónica Gil-Costa}
\ead{gvcosta@unsl.edu.ar}

\author[usach]{Carolina Bonacic}
\ead{carolina.bonacic@usach.cl}

\author[usach]{Mauricio Marín}
\ead{mauricio.marin@usach.cl}

\affiliation[usach]{organization={Departamento de Ingeniería Informática,
Universidad de Santiago de Chile},%
            city={Santiago},%
            country={Chile}}

\affiliation[unsl]{organization={Universidad Nacional de San Luis (UNSL)
and CONICET},%
            city={San Luis},%
            country={Argentina}}

\begin{abstract}
Two-dimensional shallow-water solvers underpin flood hazard mapping, but
open-source production kernels in FORTRAN or C++ offer no ergonomic path
to automatic differentiation, even as differentiable modelling unifies
geoscientific inverse problems. We present \emph{hydroflux}, a 2D
shallow-water finite-volume solver written in Rust and generic over the
numeric type: the identical code path evaluates in \texttt{f64} for
production and in forward-mode dual numbers for gradient extraction. The
scheme couples an HLLC Riemann solver, Audusse hydrostatic
reconstruction for well-balancedness, MUSCL slope-limited reconstruction
on $(\eta, u, v)$, an SSP Runge--Kutta time integrator, a point-implicit
Manning step, and a Liang--Marche flux rescaling for strictly
mass-conservative wetting and drying. We verify against a benchmark
hierarchy --- lake-at-rest preserved to machine precision
($\|\eta-\eta_0\|_\infty \approx 3 \times 10^{-16}$), the Thacker
oscillating paraboloid (relative $L^2$ 0.068\%, mass
$2 \times 10^{-5}$), the Stoker/Ritter dam-break ($L^1$ 1.0\%), a
radial dam-break, steady Manning uniform flow, and the six UK
Environment Agency 2D benchmarks --- and a matched head-to-head against
ANUGA on Stoker recovers the same accuracy class. We apply the solver
to the 2017 Aluvión Atacama event on the Río Huasco --- a semiarid
Andean reach with a 92-year DGA record --- at 30~m, with a spatially
variable Manning field derived from ESA WorldCover. As a one-day-peak
sensitivity demonstration (not a calibrated hindcast), the variable
field retains $\sim$25\% more water in the reach than a uniform value.
Forward-mode automatic differentiation through the full 2D solver
costs $1.98\times$ the \texttt{f64} wall-clock; on a matched Stoker
dam-break hydroflux integrates at $6.5\times$ the wall-clock speed of
ANUGA per simulated second. The solver, its test suite, and the
application scripts are released open-source as the verified 2D
foundation of a multi-year programme toward differentiable, GPU-native,
coupled hydrometeorological hazard simulation.
\end{abstract}

\begin{highlights}
\item A 2D shallow-water solver in Rust, generic over the numeric type for f64 + AD.
\item HLLC + Audusse + MUSCL + SSP-RK2; well-balanced, mass-conservative wet/dry.
\item Verified on Thacker, Stoker, MacDonald, UK~EA $\times 6$; head-to-head against ANUGA.
\item Forward-mode AD costs $1.98\times$ f64; $6.5\times$ faster wall-clock than ANUGA on Stoker.
\item Applied to 2017 Aluvión Atacama (Río Huasco) with ESA WorldCover Manning.
\end{highlights}

\begin{keyword}
shallow-water equations \sep flood modelling \sep finite volume \sep
automatic differentiation \sep Rust \sep well-balanced schemes
\end{keyword}

\end{frontmatter}

\linenumbers

\section{Introduction}

Two-dimensional shallow-water (SW) solvers are the computational
workhorse of flood hazard mapping, dam-break analysis, and increasingly
of compound-hazard assessment where inundation couples to sediment,
debris, and slope processes. The regulatory anchor of the field, HEC-RAS
\citep{Brunner2020}, pairs a mature FORTRAN kernel with a Windows GUI
and proprietary project files; the open-source alternatives ---
LISFLOOD-FP \citep{Bates2010, BatesDeRoo2000}, BASEMENT
\citep{Vetsch2020}, TELEMAC-MASCARET \citep{Hervouet2007}, ANUGA
\citep{Roberts2015}, Iber \citep{Blade2014}, SRH-2D \citep{Lai2010},
Delft3D \citep{Lesser2004}, and GeoClaw \citep{LeVeque2011} --- span the
regulatory and research tracks but are, almost without exception,
written in FORTRAN or C++, and none of them ships with automatic
differentiation (AD). Mature AD toolchains for compiled languages do
exist: operator-overloading libraries such as ADOL-C
\citep{Griewank1996ADOLC}, Sacado \citep{PhippsPawlowski2012Sacado}, and
CoDiPack \citep{Sagebaum2019CoDiPack} for C++; source transformation via
Tapenade \citep{HascoetPascual2013Tapenade} and TAF, which powers the
production adjoint of the MITgcm ocean model
\citep{Heimbach2005MITgcmAdjoint}; and LLVM-level differentiation via
Enzyme \citep{MosesChuravy2020Enzyme}. But retrofitting any of them onto
a hand-optimised legacy flood kernel means invasive type or build-system
surgery through code never designed for it --- which is why, three
decades into that lineage, no community shallow-water solver offers
gradients today.

Differentiable modelling has emerged over the last few years as a
unifying inverse-problem framework across the geosciences
\citep{Shen2023}, with hydrology an early adopter: differentiable,
learnable process-based models now approach state-of-the-art streamflow
accuracy \citep{Feng2022} and the ``calibration to parameter learning''
shift harnesses big-data scaling in distributed geoscientific models
\citep{Tsai2021}. In the fluid-dynamics core, JAX-Fluids
\citep{Bezgin2023} demonstrates a fully-differentiable high-order CFD
solver in JAX, and on the multi-hazard side SynxFlow \citep{Xia2025}
couples flood, landslide, and debris flow in a single GPU-accelerated
engine. The paradigm these works share is clear: the \emph{forward
solver} becomes differentiable (or GPU-native, or both), and inverse
problems --- parameter estimation, bathymetry inversion, learned
closures --- inherit efficient gradients. The differentiable layer
typically sits in a host language with tracing-based AD (JAX, PyTorch,
Julia). What remains comparatively under-served is a 2D shallow-water
flood solver that is differentiable \emph{by construction} --- designed
from the first commit around a generic numeric type rather than
retrofitted onto a legacy kernel --- in a memory-safe compiled language,
with the gradients themselves verified as part of the test suite,
exercised on the standard community benchmarks, and applied to a real
data-sparse basin. That combination --- not differentiation of compiled
code per se, which the ADOL-C-to-Enzyme lineage established --- is the
niche this paper occupies.

This paper does not claim to open that frontier; it claims to
\textbf{deliver and verify} a solver positioned within it, with two
specific design commitments. First, \emph{differentiability by numeric
genericity}: hydroflux is generic over a \texttt{Real} trait, so the
identical source compiles to \texttt{f64} for production and to a
forward-mode \texttt{Dual} type for gradient extraction. This is the
operator-overloading idiom of the ADOL-C/Sacado/CoDiPack family, applied
as a design-time commitment rather than a retrofit: every module of the
solver dispatches over the generic type, there is no tracer overhead, no
taping, and no separate adjoint code to maintain --- and, distinctively,
the gradients themselves are locked by an AD-versus-finite-differences
test suite (§2.5). Second, \emph{GIS-native verification on data-sparse
basins}: the solver ingests DEM and land-cover GeoTIFFs directly and is
exercised against the full UK Environment Agency 2D benchmark suite
\citep{NeelzPender2013} plus analytical references, then applied to a
real Chilean Andean reach with public gauge data. The contribution is
therefore a \emph{verified artifact} --- the 2D foundation of a
multi-year programme --- rather than a novelty claim about
differentiable hydraulics per se.

The paper proceeds with the governing equations and numerical scheme
(§2), a verification hierarchy from machine-precision well-balancedness
through analytical and community benchmarks (§3), an application to a
Río Huasco reach under a regulated 2017 release, with land-cover-derived
roughness (§4), the roadmap toward coupling, GPU acceleration, and
native autodifferentiation (§5), and conclusions (§6).

\section{Numerical methods}

\subsection{Governing equations}

The 2D shallow-water equations in conservative form, with bed-slope and
friction source terms:

\[\partial_t h + \partial_x (hu) + \partial_y (hv) = 0\]

\[\partial_t (hu) + \partial_x\!\left(hu^2 + \tfrac{1}{2}g h^2\right) + \partial_y (huv) = -g h\, \partial_x z_b - g h S_{fx}\]

\[\partial_t (hv) + \partial_x (huv) + \partial_y\!\left(hv^2 + \tfrac{1}{2}g h^2\right) = -g h\, \partial_y z_b - g h S_{fy}\]

with depth \(h\), velocities \((u, v)\), bed elevation \(z_b(x, y)\),
gravity \(g\), and Manning friction slopes \(S_{fx}, S_{fy}\). The state
is stored as \texttt{Conserved2D} = (h, hu, hv) on a structured
Cartesian mesh of cells indexed \((i, j)\) with row \(i\) along \(y\)
and column \(j\) along \(x\), matching the GeoTIFF raster convention so
that DEM and output rasters share a geotransform.

\subsection{Finite-volume
discretisation}

Cell-centred finite volume with the HLLC approximate Riemann solver
\citep{Toro1994, Toro2009} evaluated per face via the
rotational-invariance of the SW flux (a 1D normal-direction Riemann
problem at each x- and y-face). The bed-slope source is treated in the
well-balanced hydrostatic-reconstruction framework of Audusse et al.
\citep{Audusse2004}: at each face the reconstructed depths

\[h^*_L = \max(h_L + z_L - z_{\max}, 0), \quad h^*_R = \max(h_R + z_R - z_{\max}, 0)\]

with \(z_max = max(z_L, z_R)\) feed the HLLC flux, and the
hydrostatic-pressure correction \((g/2)(h^{2} - h*^{2})\) is added to
the face-normal momentum component on each side. The bed is
reconstructed linearly to the face midpoint following Liang \& Marche
\citep{LiangMarche2009}, which moves the bed-slope source from the face
flux into a cell-centred algebraic term
\(S = (g/2)(h^{2}_{R,\text{face}} - h^{2}_{L,\text{face}})/\Delta x\)
that cancels the pressure-flux divergence exactly at rest (the
C-property --- exact preservation of a flat free surface over an
arbitrary bed; verified to machine precision in §3.1).

\subsection{Second-order reconstruction and time
integration}

Second-order spatial accuracy uses MUSCL reconstruction with the minmod
limiter applied to the \textbf{primitive} vector \((\eta , u, v)\),
where \(\eta = h + z_b\) is the free-surface elevation --- not the
conserved variables and not \((\eta , hu, hv)\). The choice follows the
well-balancing argument of Liang and Marche \citep{LiangMarche2009}:
reconstructing velocities rather than momenta is what makes the scheme
well-balanced on flows with non-uniform depth over a sloped bed. When
the analytical solution has uniform \(u\) but varying \(h\)
(e.g.~Manning normal flow), \((\eta , u, v)\)-MUSCL gives equal velocity
on both sides of every face and the HLLC sees a consistent state,
eliminating an \(O(\Delta x\, S_0/h_n)\) steady-state drift that
momentum-reconstruction incurs (§3.4 measures a 45× reduction versus
\((\eta , hu, hv)\)-MUSCL on the steady Manning uniform-flow test).

Time integration uses the strong-stability-preserving Runge--Kutta
second-order method (SSP-RK2; Heun), written as the convex combination
of two forward-Euler updates so that every property the forward-Euler
step preserves under the CFL bound --- non-negative depth, mass
conservation, finite momentum --- is inherited \citep{Toro2009}. The
time step is bounded by
\(dt \cdot (s_x/\Delta x + s_y/\Delta y) \leq \text{CFL}\) with
\(s_x = \max(|u| + c)\), \(s_y = \max(|v| + c)\), \(c = \sqrt{g h}\).

\subsection{Friction, wetting/drying, and boundary
conditions}

Manning friction is applied as an operator-split point-implicit
fractional step on the momentum,
\((hu, hv) \leftarrow (hu, hv)/(1 + \beta)\) with
\(\beta = \Delta t\, g\, n^2 |U| / h^{4/3}\), which is unconditionally
stable, preserves rest and dry cells exactly, and keeps the flow
direction fixed (the shared factor divides both momentum components).
The Manning coefficient \(n\) is stored as a per-cell field, allowing a
spatially variable roughness derived from land cover (§4.2); a uniform
scalar is the special case of a constant field.

Wetting and drying use the Liang \& Marche flux-rescaling scheme: a
per-cell factor \(\alpha \in [0, 1]\) caps the total outgoing mass at
the available depth, applied to all three flux components at each
outgoing face, with the upstream cell's factor seen by both cells
sharing the face. This guarantees no cell drains below the dry threshold
and that mass is conserved exactly across each face; on wet--wet flows
\(\alpha ≡ 1\) and the scheme reduces to the unrescaled finite-volume
update. Boundary conditions (Wall, Transmissive, Discharge, Depth) are
imposed via ghost cells on all four sides; the Discharge-on-dry case
reconstructs a Manning-normal-depth ghost so a cold-start inflow over a
dry domain picks a sane first time step.

\subsection{Differentiability and
performance}

The 2D solver is generic over a \texttt{Real} trait abstracting the
arithmetic surface (\(+\), \(-\), \(×\), \(/\), \texttt{sqrt},
\texttt{powf}, \texttt{max}, \texttt{min}, \texttt{abs}). Two
implementations are used: \texttt{f64} for production, and a
forward-mode \texttt{Dual\ \{val,\ dval\}} propagating
\((value, derivative)\) through every operation. The state
(\texttt{Conserved2DG\textless{}T\textgreater{}}), the mesh
(\texttt{Mesh2DG\textless{}T\textgreater{}}), the HLLC Riemann flux, the
Audusse hydrostatic reconstruction, the MUSCL slopes, the Liang--Marche
flux rescaling, the point-implicit Manning step, and both time
integrators (\texttt{forward\_euler\_step}, \texttt{ssprk2\_step}) all
dispatch over \(T: Real\). The derivative of any output with respect to
a single seed input is recovered as \texttt{result.dval} after one
forward pass; there is no separate adjoint code to maintain. Branching
on dryness uses \(T::value()\) so control flow stays scalar; on each
branch the arithmetic propagates the derivative through whichever side
is taken.

We verify this end-to-end with an AD-vs-FD locking suite that compares
forward-mode AD against second-order central finite differences on
independent gradients across every layer of the discretisation. On the
HLLC Riemann flux (wet/wet star branch and the dry-bed two-rarefaction
branch), the Manning friction step (seeded on both \(n\) and \(h\)), the
full forward-Euler update, and the SSP-RK2 update, AD matches central FD
to better than \(10^{-6}\) relative on derivatives whose magnitudes span
six orders of magnitude. The SSP-RK2 test additionally recovers an exact
analytical invariance: the derivative of total mass with respect to a
uniform bed-elevation shift is
\(-(n_rows \cdot n_cols \cdot dx \cdot dy)\) to \(10^{-8}\) --- mass
conservation under reference-level change holds in the gradient, not
just the value.

Two classes of non-smoothness require an explicit choice, and we state
both rather than leave them to the arithmetic. The first is the kinks
that wetting and drying introduce into the flux. At \(h = 0\) the
derivative of \texttt{sqrt} is unbounded, so the clamp-to-dry
composition \(\sqrt{\max(h, 0)}\) would return NaN at the shoreline; we
return the subdifferential element \(0\) instead, which keeps the
gradient finite and zero exactly where the cell carries no water.
\texttt{abs} at the origin and \texttt{max}/\texttt{min} at a tie are
resolved the same way --- \(0\) for the former, the average of the two
incoming derivatives for the latter --- so the gradient stays bounded
and symmetric when an optimiser lands on the kink. The second is the
dryness thresholds \texttt{H\_DRY} and \texttt{H\_VEL}: these branch on
\(T::value()\), so control flow is scalar and the derivative propagates
through whichever branch the primal takes. The resulting map is
piecewise differentiable, and the gradient is a one-sided derivative of
the active piece rather than of a smoothed surrogate.

The time step is the one deliberate omission. \texttt{cfl\_dt} extracts
\(.value()\) and returns \texttt{f64}, so the CFL sequence is treated as
a constant schedule and is not differentiated. The recovered gradient is
therefore that of the scheme with a \emph{frozen} \(dt\) sequence, which
differs from the true sensitivity of the discrete map by an \(O(dt)\)
term --- the semi-implicit friction factor depends on \(dt\) through
\(|q + dt\cdot R|\), so \(\partial G/\partial dt \neq 0\) even at a
steady fixed point. We measure that gap rather than assume it is
negligible: on the 1D power-law steady-state test it is
\(1.2\cdot 10^{-3}\), \(2.2\cdot 10^{-3}\) and \(1.7\cdot 10^{-3}\)
relative for the three calibrated parameters at CFL 0.4, and it halves
when the CFL number halves --- the signature of an \(O(dt)\) term rather
than a defect in the dual-number rules, which would appear at \(O(1)\).
For transient quantities of interest the gap is larger and the tolerance
should be re-derived; for calibration against steady or slowly varying
targets it sits three orders of magnitude below the parameter
sensitivities themselves.

A second and more consequential limit is the integration horizon, and we
state it because the locking suite above does not reveal it. Those tests
integrate O(100) steps on smooth synthetic problems. Repeating the same
seeded evaluation on the §4 configuration --- a sealed 30 m reach with a
moving shoreline, $\sim$78 000 steps for one simulated day ---
the tangent grows exponentially at 1.85 \% per step --- a log-linear fit
over 51 sampled steps gives \(0.0080 \pm 0.0002\) decades per step
(1\(\sigma\)), 95 \% interval 1.78--1.93 \% per step, \(R^{2} = 0.98\),
and the quality of that fit is itself the evidence that the growth is
exponential rather than merely large --- while the primal remains
entirely stable: peak depth drifts from 2.40 m to 2.94 m and the
wet-cell count from 219 to 220 over the same interval. The consequence
is stark. Over 100 steps the amplification is a factor of six and
invisible; over a full day it is \(10^{606}\), and the gradient
overflows. A Gauss-Newton calibration of the §4.4 roughness target on
this horizon therefore returns a meaningless step, which is what we
observe.

This is a property of the linearisation, not of the differentiation
mode: an adjoint of the same trajectory inherits the same instability,
so reverse-mode AD would not repair it. What it bounds is the class of
inverse problem this solver currently supports --- short assimilation
windows, or steady and near-steady targets, where the tangent stays
bounded --- and not long transient hindcasts. Methods that address
unstable tangents directly, such as shadowing or windowed
regularisation, are the honest route and are outside this paper's scope.
We report the measurement rather than the ambition:
\texttt{diag\_dual\_growth} reproduces the growth curve.

Beyond locking, we exercise the wedge on a small inverse problem. A
friction-damped dam-break on a 3×60 channel with
\((h_L, h_R) = (1.0, 0.1)\) m is simulated to \(t = 20\) s; the
depth-weighted centroid of the wet region serves as the observation.
Picking \(n_true = 0.040\) produces a synthetic centroid at
\(x = 27.65\) m. Starting from \(n_init = 0.080\) (a 2× over-estimate)
and updating
\(n \leftarrow n - (c_pred(n) - c_obs) \cdot (\partial c_pred/\partial n)^{-1}\)
with the gradient from one forward pass of \(T = Dual\), Newton
converges to \(n_final = 0.040 000\) (relative error \(< 10^{-15}\)) in
five iterations.

The arithmetic overhead of the wedge is measured on a 64×64 cell domain
over 200 SSP-RK2 + Manning steps. The Dual-typed run takes 2.01× the
wall-clock of the \texttt{f64} run (\(549\) vs \(273\) ns per cell-step,
post-warm-up, release build). This is within the 2-3× band expected for
operator-overloading forward AD on compiled code
\citep{Griewank2008, Sagebaum2019CoDiPack}.

The \texttt{Dual} type carries one derivative direction, so a gradient
with respect to \(P\) parameters costs \(P\) independent forward passes.
We measure that scaling rather than estimate it: sweeping
\(P \in \{1, 2, 4, 8, 16\}\) zonal Manning coefficients on a 64×64
dam-break, the cost of the full gradient is linear in \(P\) at
\(2.07\times\) the \texttt{f64} primal per parameter (per-parameter
ratios 2.04, 2.04, 2.04, 2.06, 2.16), and the \(P = 1\) point recovers
the \(2.0\times\) single-seed overhead reported above. Reverse-mode AD
recovers the whole gradient in one sweep at a cost independent of \(P\),
conventionally 3--5× the primal, which puts the break-even at
\(P^{*} \approx 2\) (1.5--2.4 across that band). Forward mode is
therefore the right tool only for the low-dimensional targets this paper
exercises --- and §4.4 shows the roughness problem here \emph{is}
low-dimensional, with a single land-cover class carrying almost all of
the sensitivity. It is not the right tool for a per-cell roughness
field, and we do not claim otherwise; reverse-mode is identified as
future work (§5). The counterpoint that does favour forward mode is
memory: it stores no tape, so its footprint is independent of the number
of time steps, whereas an adjoint must checkpoint or replay a trajectory
that in these simulations runs to \(10^{5}\) steps. The 1D companion
line \citep{ParraPaper02} uses the same \texttt{Real} trait and the same
\texttt{Dual} type, exposed through the autograd crate. Reproduced by
the \texttt{m1\_forward\_scaling} example.

A cell-mask early-skip optimisation exploits the fact that arid-basin
simulations are dominated by dry cells: interior faces with both
adjacent cells dry return zero flux without evaluating the HLLC or the
MUSCL reconstruction (the dry--dry flux is identically zero, not merely
small, so the skip is exact), and strictly-interior cells with all four
neighbours dry skip the cell update. On the Huasco application (§4),
where $\sim$97 \% of cells are dry, this cuts the wall time by
$\sim$38 \% (a 1.6× speed-up) with bit-identical results. The
solver is \texttt{\#!{[}forbid(unsafe\_code){]}}.

\section{Verification}

We verify against a hierarchy of references of increasing complexity,
from exact steady states through analytical transients to the community
benchmark suite. Table 1 summarises the quantitative results; all are
computed by the solver's automated test suite and reproduced by the
\(report_*\) informational tests.

A word on what the uncertainties in this paper are, since the three
classes of number reported here are not alike. The verification metrics
of Table 1 carry no run-to-run uncertainty at all: the solver contains
no stochastic element and no parallel reduction whose order can vary, so
repeated invocations return bit-identical values, which we confirmed by
running the Thacker report three times. Their uncertainty is entirely
discretisation, and §3.7 characterises it directly by refining the mesh.
The performance figures of §3.9 are the opposite case --- machine-
dependent and subject to run-to-run scatter, which is why they are
reported with the hardware stated and, where measured on a different
machine from the rest, flagged as such. The third class, the inter-model
comparison of §4.5, has an uncertainty that neither of these covers, and
§4.5 quantifies it.

\textbf{Table 1. Verification results.}

\begin{longtable}[]{@{}
  >{\raggedright\arraybackslash}p{(\columnwidth - 8\tabcolsep) * \real{0.2821}}
  >{\raggedright\arraybackslash}p{(\columnwidth - 8\tabcolsep) * \real{0.1538}}
  >{\raggedright\arraybackslash}p{(\columnwidth - 8\tabcolsep) * \real{0.1538}}
  >{\raggedright\arraybackslash}p{(\columnwidth - 8\tabcolsep) * \real{0.2051}}
  >{\raggedright\arraybackslash}p{(\columnwidth - 8\tabcolsep) * \real{0.2051}}@{}}
\toprule\noalign{}
\begin{minipage}[b]{\linewidth}\raggedright
Benchmark
\end{minipage} & \begin{minipage}[b]{\linewidth}\raggedright
Type
\end{minipage} & \begin{minipage}[b]{\linewidth}\raggedright
Mesh
\end{minipage} & \begin{minipage}[b]{\linewidth}\raggedright
Metric
\end{minipage} & \begin{minipage}[b]{\linewidth}\raggedright
Result
\end{minipage} \\
\midrule\noalign{}
\endhead
\bottomrule\noalign{}
\endlastfoot
Lake-at-rest, bumpy bed & analytical (C-property) & 20×20 &
\(\|\eta - \eta _{0}\|_\infty\) & \(< 10^{-10}\) (test bound) \\
Lake-at-rest, Thacker paraboloid & analytical (C-property) & --- &
\(\|\eta - \eta _{0}\|_\infty\) & \(\approx 3\cdot 10^{-16}\)
(measured) \\
Thacker oscillating & analytical transient & 80×80 & rel. L$^{2}$ on \(h\) &
0.0734 \% \\
Thacker oscillating & analytical transient & 80×80 & mass error &
\(3.53\cdot 10^{-15}\) \\
Stoker/Ritter dam-break & analytical transient & 400×3 & L$^{1}$ on \(h\) &
1.0 \% \\
Stoker/Ritter dam-break & analytical transient & 400×3 & L∞ on \(h\) &
2.2 \% \\
Radial dam-break & symmetry & 160×160 & axisymmetry & preserved \\
MacDonald uniform flow & steady, well-balanced & 5×50 & steady-state
\(h\) & 0.073 \% (guard 2 \%) \\
UK EA Tests 1--6 (synthetic stand-ins) & community suite & various &
qualitative + mass & pass \\
UK EA Test 4 (official EA/LISFLOOD-FP geometry) & community suite,
quantitative & 400×200 @ 5 m & RMSE vs LISFLOOD-FP DG2 (6 control
points) & 0.3--1.2 \% of peak depth \\
UK EA Test 8A (official EA/LISFLOOD-FP geometry) & community suite,
qualitative & 481×199 @ 2 m & vs SC120002 inter-model bounds (4 scored
points) & pass, incl.~inside industry spread \\
\end{longtable}

\subsection{Well-balancedness
(lake-at-rest)}

A flat free surface over an arbitrary bed must remain at rest. On a
submerged 2D Gaussian bump and on a smooth Thacker paraboloid, the
scheme preserves \(\eta = h + z_b\) and zero momentum to machine
precision (\(\|\eta - \eta _{0}\|_\infty \approx 3\cdot 10^{-16}\),
\(\|q\|_\infty \approx 2\cdot 10^{-15}\) after 60 s). The cell-centred
algebraic bed-slope source of §2.2 is bit-exact for \emph{any} bed shape
--- including smooth curved beds --- provided every cell stays wet,
correcting an earlier working assumption in our own development that
smooth curved beds would require a Castro--Parés path-conservative
treatment. The cancellation is self-consistent in the face beds
\texttt{z\_face}, not in the underlying bed function.

\subsection{Thacker oscillating
paraboloid}

The Thacker planar-oscillation solution on a paraboloidal basin
\citep{Thacker1981} is a classic 2D analytical transient with a moving
wet/dry shoreline. Over a half-period (388 steps on an 80×80 mesh), the
solver reproduces the analytical depth with relative L$^{2}$ error 0.0734 \%
and L∞ error \(2\cdot 10^{-4}\) m (0.17 \% of \(h_{0}\)), conserving
mass to \(3.53\cdot 10^{-15}\) despite the continuously moving shoreline
--- ten orders of magnitude tighter than an earlier measurement on this
benchmark, consistent with the wet/dry threshold treatment of §2.4 (a
thin residual film is kept rather than zeroed at the shoreline, so mass
is not discarded as the front sweeps back and forth).

\subsection{Dam-break on a dry bed
(Stoker/Ritter)}

The Ritter/Stoker dam-break \citep{Stoker1957} has a closed-form
rarefaction-plus-front solution. On a 400-cell channel with \(h_L = 1\)
m, the SSP-RK2 scheme attains L$^{1}$ error 1.0 \%, L$^{2}$ 1.0 \%, and L∞ 2.2 \%
of \(h_L\) at \(t = 4\) s, with the wet front lagging the analytical
position by 3.18 m (the expected diffusive lag of a shock-capturing
scheme at the dry front). The forward-Euler integrator gives a
comparable L$^{1}$ (1.1 \%) and an identical front lag (3.18 m), confirming
that the spatial scheme, not the time integrator, dominates the error
budget.

\subsection{Steady Manning uniform
flow}

We use the degenerate limit of the MacDonald inverse-design family
\citep{MacDonald1997}: steady Manning uniform flow at constant normal
depth \texttt{h\_n} on a uniformly sloped bed, where the bed-slope
gravity term is balanced exactly by Manning friction. The target profile
is flat, but the test is non-trivial in what it exercises simultaneously
--- the well-balanced bed-slope source, the point-implicit friction
step, and the Discharge (upstream) and Depth (downstream) boundary
conditions --- and it is the configuration on which momentum-vector
reconstruction visibly fails. The solver holds the prescribed
\texttt{h\_n} to 0.073 \% and the prescribed unit discharge to 0.18 \%
(against a 2 \% regression guard), more than an order of magnitude
better than a \((\eta , hu, hv)\)-momentum reconstruction, which
motivated the primitive \((\eta , u, v)\) choice of §2.3.

\subsection{Radial dam-break and
isotropy}

A circular dam-break on a 160×160 mesh tests grid isotropy: the radial
depth profile must be independent of azimuth. The solver preserves
axisymmetry --- depths along \(+x\) and \(+y\) are bit-identical, and
the diagonal (\(+45°\)) profile agrees to $\sim$1 \% ---
confirming the x/y flux assembly carries no directional bias.

\subsection{UK Environment Agency 2D benchmark
suite}

The UK EA 2D benchmark report \citep{NeelzPender2013} specifies ten
configurations (Tests 1, 2, 3, 4, 5, 6A, 6B, 7, 8A, 8B) built to
exercise features that matter operationally in flood modelling. We use
this suite in two complementary ways: six lightweight synthetic
stand-ins as fast CI regression tests, and --- on top of that --- two of
the tests reproduced on the \emph{official} EA geometry with a
quantitative or semi-quantitative comparison against published reference
results.

\textbf{Synthetic stand-ins (Tests 1--6,
\texttt{solver-2d/tests/uk\_ea\_test*.rs}).} These capture the essential
physics of the analogous official configurations --- filling of a
disconnected low-lying pond, rainfall on a floodplain, flow past an
obstruction, long-wave propagation in a valley, valley flooding with a
parabolic cross-section, and an urban dam-break through a building array
--- without reproducing the exact EA geometry files. The solver passes
all six, with buildings remaining dry and no spurious oscillations at
the wet/dry fronts (the strict mass-conservation figures are reported
for the closed-domain tests in §3.1--§3.2; the UK EA cases use open
inflow/outflow boundaries). These run in seconds and stay in CI as
regression guards; they are not presented as a validation against the
official benchmark.

\textbf{Official-geometry reproduction (Tests 4 and 8A).} For a
quantitative claim, we reproduce two of the ten official configurations
directly from EA/LISFLOOD-FP data redistributed under CC-BY-4.0 by the
LISFLOOD-FP team \citep{Shaw2021, Sharifian2023} --- the proprietary EA
originals are not publicly downloadable.

\emph{Test 4} (``speed of propagation of a flood wave'') uses the
official 1000 m × 2000 m flat floodplain (5 m mesh, Manning \(n=0.05\)),
the official trapezoidal hydrograph peaking at 20 m$^{3}$/s through a 20 m
inlet, and the six official control points. Comparing simulated depth
against the redistributed LISFLOOD-FP DG2 reference series (5 m, full
second-order shallow-water scheme --- the resolution- and scheme-matched
comparison) gives RMSE of 0.3--1.2 \% of peak depth at every point, peak
bias within ±1.4 \%, and arrival-time offsets of 0--60 s --- an order of
magnitude below the $\sim$5 min spread the SC120002 report
itself documents \emph{between different industry models} on this test.
Full detail in
\texttt{benchmarks/data/uk\_ea/test4/results\_hydroflux.md}.

\emph{Test 8A} (``rainfall and point source surface flow in urban
areas'', Cockenzie Street, Glasgow) uses the official 2 m DEM, a
spatially variable Manning field (roads vs.~background), a 400 mm/h
3-minute rainfall pulse, and a point inflow peaking at 5 m$^{3}$/s. No
numeric LISFLOOD-FP reference series is redistributed for this test, so
the comparison is against the \emph{qualitative bounds reported in the
text} of SC120002 §4.9.3 --- agreement ranges observed across the
$\sim$15 industry packages that ran the official test. All four
control points with a stated numeric bound pass, including the
downstream pond (point 3): simulated final depth is within 0.066 m of
the expected $\sim$0.8 m, inside the $\sim$0.07 m
spread the report documents \emph{between} the $\sim$15
industry models. Full detail in
\texttt{benchmarks/data/uk\_ea/test8a\_glasgow/results\_hydroflux.md}.

We attempted a similar official-geometry reproduction of Test 5
(``valley flooding''); the official DEM is not in any public
redistribution (only proprietary from the EA), and an idealised
synthetic valley built from the report's text description alone produced
order-of-magnitude but not quantitatively meaningful agreement,
dominated by the necessarily invented inflow hydrograph rather than by
solver behaviour --- we do not include it as evidence here.

\subsection{Mesh-refinement
convergence}

We quantify the order of accuracy on the Thacker oscillation --- the
benchmark here that combines a smooth analytical transient with a moving
wet/dry shoreline --- by refining the grid from 32$^{2}$ to 256$^{2}$ and
measuring the relative L1 and L2 error in depth at \(t = T/2\) (Table 2,
Fig.~\ref{fig:convergence}).

\textbf{Table 2. Thacker convergence (error in \(h\) at \(t = T/2\)).}

\begin{longtable}[]{@{}llllll@{}}
\toprule\noalign{}
\(n\) & \(\Delta x\) {[}m{]} & rel. L1 & rel. L2 & order L1 & order
L2 \\
\midrule\noalign{}
\endhead
\bottomrule\noalign{}
\endlastfoot
32 & 0.0781 & 4.82·10$^{-3}$ & 5.78·10$^{-3}$ & --- & --- \\
64 & 0.0391 & 1.38·10$^{-3}$ & 1.76·10$^{-3}$ & 1.81 & 1.71 \\
128 & 0.0195 & 4.05·10$^{-4}$ & 5.89·10$^{-4}$ & 1.77 & 1.58 \\
256 & 0.0098 & 1.35·10$^{-4}$ & 2.18·10$^{-4}$ & 1.59 & 1.43 \\
\end{longtable}

The observed order starts near second (1.81 in L1 between the two
coarsest grids) and relaxes toward $\sim$1.5 at the finest,
with overall log-log slopes of 1.73 (L1) and 1.58 (L2). This is the
expected signature of a formally second-order scheme (MUSCL + SSP-RK2)
whose moving wet/dry shoreline is locally first order: as the
smooth-interior error shrinks under refinement, the first-order
shoreline becomes a larger share of the total and bends the global rate
below two. The lake-at-rest tests (§3.1) confirm the interior
discretisation is exact to machine precision, so the sub-two rate
reflects the wet/dry front treatment rather than a deficiency in the
smooth-region scheme --- consistent with the orders reported for
comparable well-balanced shallow-water solvers \citep{LiangMarche2009}.

\subsection{Head-to-head against
ANUGA}

To position the solver against a mature open-source 2D shallow-water
reference, we ran the Stoker dam-break in ANUGA \citep{Roberts2015} ---
its default DE0 flow algorithm on the \texttt{rectangular\_cross}
triangulation, built on the central-upwind scheme of Kurganov and
Petrova \citep{KurganovPetrova2007} --- at the same effective resolution
(\(\Delta x = 1\) m, matched to a 100-cell hydroflux re-run) and the
same physical setup (flat bed, walls on the long sides, transmissive
ends, \(h_L = 1\) m, \(t_end = 4\) s). The two schemes are of different
families --- a Riemann-solver discretisation here against a Riemann-free
central-upwind one there --- which is what makes agreement between them
informative rather than circular. Both reproduce the Ritter rarefaction
solution closely (Fig.~\ref{fig:head_to_head}). On the analytical rarefaction fan
\(x ∈ [37.5, 75.1]\) m, the relative error norms are L1 4.1 \%, L2 3.6
\%, L∞ 5.3 \% \(h_L\) for hydroflux and L1 2.6 \%, L2 2.7 \%, L∞ 4.4 \%
\(h_L\) for ANUGA --- ANUGA edges out hydroflux by a factor of
$\sim$1.5 in L1 at this coarse resolution, with both staying
below 6 \%. Refining hydroflux to its reference 400-cell mesh
(\(\Delta x = 0.25\) m, the configuration of Table 1 and §3.3) drops the
L1 error to 1.0 \%, consistent with the near-second-order convergence of
§3.7. The two solvers are therefore in the same accuracy class on this
canonical dry-front benchmark; the (modest) coarse-grid gap closes under
refinement, as the convergence rates predict.

\subsection{Computational
performance}

We characterise serial throughput on synthetic grids and report a
wall-clock comparison against ANUGA on the same Stoker problem the
accuracy comparison above used. Serial throughput and the AD-overhead
measurement below are from a release build on a quiet 8-physical-
core/12-thread x86-64 laptop (Intel i5-13420H, native Linux); the ANUGA
wall-clock comparison that follows was measured earlier on a different
(16-core) workstation and was not re-run here (no ANUGA installation on
the quiet machine) --- the two are not cross-comparable against each
other, only within their own pairing. The corresponding benchmark
scripts (\texttt{m2\_perf\_large\_grid.rs},
\texttt{m2\_hydroflux\_wallclock.rs}, \texttt{m2\_anuga\_wallclock.py})
are released with the code.

\textbf{Serial throughput.} On smooth Gaussian-bump initial conditions
(no dry interior --- the cell-mask early-skip optimisation cannot
artificially flatter the timing), the solver sustains \(3.4\)--\(3.6\)
Mcell-steps per second at \texttt{f64} precision across grids from
\(256^{2}\) to \(1024^{2}\): \(282\), \(292\) and \(294\) ns per
cell-step respectively at the three sizes. The throughput is consistent
with the single-threaded CPU references cited by the
Caviedes-Voullième-circle GPU papers \citep{Rak2024, SaleemNorman2024}
and falls in the band typical of HLLC + MUSCL + SSP-RK2 implementations
on cache-resident problems. Community GPU/HPC shallow-water solvers ---
SERGHEI-SWE \citep{CaviedesVoullieme2023SERGHEI}, scaled to hundreds of
GPUs on TOP500-class systems, and the multi-GPU TRITON
\citep{MoralesHernandez2021TRITON} --- demonstrate the throughput
headroom a SIMT port unlocks over single-threaded CPU execution; §5(i)
motivates the \(wgpu\) port this comparison sets up.

\textbf{\texttt{f64} vs \texttt{Dual\textless{}f64\textgreater{}}
overhead.} On a \(64^{2}\) grid over \(200\) SSP-RK2 + Manning steps,
the forward-mode AD instance takes \(2.01 ×\) the wall-clock of the
\texttt{f64} instance (\(273\) vs \(549\) ns per cell-step, §2.5). The
ratio is inside the 2-3× band expected for operator-overloading forward
AD on compiled code \citep{Griewank2008, Sagebaum2019CoDiPack}. We treat
this as the AD overhead headline.

\textbf{Wall-clock vs ANUGA on a matched Stoker problem.} On the Stoker
dam-break scaled to \(200 m × 5 m\), \(t_end = 8 s\) at matched
effective \(\Delta x = 0.5 m\), hydroflux integrates the eight simulated
seconds in \(1.07 s\) of wall-clock (\(0.13\) s wall per simulated
second); ANUGA integrates the same physical problem in \(6.91 s\) of
wall-clock (\(0.86\) s wall per simulated second). The ratio is \(6.5×\)
in hydroflux's favour per simulated second. The two solvers run
different mesh topologies (rectangular cells versus the
\texttt{rectangular\_cross}-triangulated \(4 ×\) denser mesh that
ANUGA's solver requires) so cell-count throughput is not directly
comparable; the simulated-second metric is the meaningful one because
the user- visible cost of running a flood event scales with this
quantity.

\textbf{Bottleneck and CPU parallelism.} Profiling the forward Euler
step (rough profile from \texttt{cargo\ flamegraph} on the \(512^{2}\)
Gaussian bump) shows that the \texttt{well\_balanced\_x\_face} +
\texttt{well\_balanced\_y\_face} assembly takes $\sim$60 \% of
the wall-clock, the cell update (with the explicit bed-slope source)
$\sim$25 \%, and the cell-mask + MUSCL slope pass the
remainder. Both dominant pieces are embarrassingly parallel on a
per-face/per-cell basis. An initial attempt at CPU parallelism, one
\(rayon\) task per face, was defeated by task-dispatch overhead relative
to the $\sim$200-500 ns of per-face arithmetic. Re-measured at
\textbf{row granularity} instead --- \texttt{ndarray}'s
\texttt{Zip::par\_for\_each} over each \texttt{StepWorkspace2D} pass,
which splits work by recursively bisecting the outer axis rather than
per-element --- CPU parallelism \emph{is} worthwhile: on a \(256^{2}\)
all-wet grid (quiet 8-core/12-thread machine, release build), the
reusable-workspace step reaches \(3.2×\) at 4 threads and saturates at
\(3.8\)--\(4.0×\) by 8 threads (\texttt{ssprk2\_step\_with}: serial 13.9
ms → 3.6 ms). On a mostly-dry grid resembling the §4 application
($\sim$94 \% dry, dominated by the wet/dry short-circuit), the
same measure saturates lower, at \(2.8×\): less per-cell work survives
the dry skip to parallelise, and row chunks are unevenly loaded when the
wet channel occupies a small, uneven fraction of rows. Both regimes
plateau by 4-8 threads; the remaining threads on our 12-thread test
machine (8 physical cores) buy nothing further. This is now a
feature-gated capability (\texttt{-\/-features\ parallel},
\texttt{T:\ Send\ +\ Sync} required only on that path, not on the
\texttt{Real} trait), not a research dead end --- but its ceiling is
well below the throughput SIMT execution is expected to offer on the
same embarrassingly parallel face/cell loops. Section §5 carries the GPU
port (via \(wgpu\) compute shaders) as the immediate next deliverable.

\section{Application: a Río Huasco reach under a regulated 2017
release}

\subsection{Setup}

The Río Huasco drains a semiarid Andean basin in the Atacama region of
northern Chile, a setting whose episodic flow regime is well documented
\citep{Wilcox2016AtacamaFlash, Cabre2020HuascoENSO}. We model a 200 ×
67-cell reach of the 30 m pit-filled SRTM DEM (6 km × 2 km, UTM 19S, bed
elevations 461--888 m), ingested directly as a GeoTIFF. Boundary
conditions are Transmissive on the western (downstream) edge and Wall on
the others, with a point-source inflow at the eastern channel cell.

The inflow is the daily discharge recorded at DGA gauge 03820003 (Río
Huasco en Santa Juana) for the 21-day window 2017-02-20 → 2017-03-12,
peaking at 38.9 m$^{3}$/s on 2017-03-02. Two properties of that record govern
how the experiment should be read, and we state both explicitly.

First, the gauge lies 3.5 km upstream of the modelled reach --- outside
the domain --- so the series is a measured upstream boundary condition
rather than an interior observation the simulation could be scored
against. Second, and more consequentially, the gauge sits below the
Santa Juana reservoir: the national water authority segments this
stretch of the Río Huasco at the reservoir's inlet and outlet, placing
the Algodones gauge on the inflow and Santa Juana on the release
\citep{DGA2004Huasco}. The 2017 series is therefore an operated release
rather than a natural flood wave, and the surrounding gauge network
makes that unambiguous: over the same window the next gauge upstream
(Chepica) rises only 11 \% (12.6 → 14.0 m$^{3}$/s) while Santa Juana rises
122 \% (17.5 → 38.9 m$^{3}$/s), a difference no natural routing over 10 km of
desert without a tributary can produce, and the release ramps over six
days and recedes over three weeks rather than showing the sharp rise and
fall of a semiarid flash flood. Consistent with this, the
Richards--Baker flashiness index at the gauge falls from 0.078 over
1928--1994 to 0.045 over 1996--2019.

We use the release deliberately rather than in spite of its origin. The
experiment of §4.3 isolates the effect of the roughness field, and a
gauged release is a \emph{better} upstream condition for that purpose
than a hydrograph inferred from rainfall--runoff: the forcing carries no
rainfall-runoff or routing uncertainty, so the difference between the
two Manning configurations is attributable to the friction field alone.
It is not uncertainty-free --- the series is a rated discharge, so it
inherits the rating-curve uncertainty of any gauged record --- but that
uncertainty enters both configurations identically and therefore cancels
in the comparison, which is what the experiment needs. What the setup
does not support is a claim of hindcast skill, and none is made. The
daily-mean series is the finest resolution the public DGA record
provides here; smoothing the sub-daily shape makes simulated peak
inundation conservative for a given daily volume, which is again
consistent with the sensitivity scope of §4.3. Channel cells are
warm-started at the Manning normal depth for the day-1 discharge.

\subsection{Spatially variable Manning from land
cover}

The solver ingests an ESA WorldCover 2021 land-cover raster
\citep{ESAWorldCover2021}, resampled from its native 10 m to the 30 m
DEM grid by majority (mode) resampling, and maps each class to a Manning
coefficient through a published lookup (Chow \citep{Chow1959} and
compilations therein): bare/sparse ground (66 \% of the domain,
\(n = 0.025\)), grassland (14 \%, \(n = 0.040\)), tree cover (8 \%,
\(n = 0.100\)), shrubland (8 \%, \(n = 0.060\)), cropland (3 \%,
\(n = 0.035\)), built-up (1 \%, \(n = 0.015\)), and permanent water
(\textless{} 0.1 \%, \(n = 0.030\)). The land cover is not random with
respect to the channel: riparian tree and shrub vegetation
(\(n = 0.06\)--\(0.10\)) tracks the thalweg, while the surrounding
hillslopes are bare desert (\(n = 0.025\)). The resulting field has
\(n_{\min} = 0.015\), \(n_{\text{mean}} = 0.036\), \(n_{\max} = 0.100\).
The mapping is illustrative rather than calibrated: the absolute values
are within the standard Manning ranges \citep{Chow1959} but a regional
friction survey or a literature compilation with narrower-than-Chow
uncertainties (e.g.~Arcement and Schneider's USGS compilation) would
tighten the lookup. Because that choice propagates directly into the
headline result of §4.3, we measure its influence with a one-at-a-time
sweep (§4.4) rather than assert it.

\subsection{Results}

We route the release from the start of the window to its peak on day 11
(38.9 m$^{3}$/s) and compare the spatially variable Manning field against a
single calibrated \(n = 0.04\). At the peak the variable field raises
the mean channel depth by 0.33 m and retains 29 \% more water in the
reach (\(3.92\cdot 10^{5}\) vs \(3.03\cdot 10^{5}\) m$^{3}$), with the
wetted-cell count 277 → 302 against 295 for the uniform field. The
mechanism is physical: the riparian vegetation that the land cover
places in the channel (\(n \approx 0.10\), four times the uniform value)
slows the flow, deepens it, and holds more water in the reach --- an
effect a single domain-averaged roughness cannot represent. Mass is
conserved throughout.

Two aspects of that comparison depend on discharge, and we report both
rather than the more flattering one. The effect on storage strengthens
as the event builds: the same comparison run at the day-one baseline of
17.5 m$^{3}$/s gives +22 \% rather than +29 \%, and a mean channel deepening
of 0.19 m rather than 0.33 m. The effect on through-flow does the
opposite and nearly vanishes: mean outflow differs by 4 \% at baseline
(15.0 vs 15.6 m$^{3}$/s) but by 0.4 \% at the peak (27.9 vs 28.0 m$^{3}$/s),
because after eleven days of routing both configurations are close to
steady state and the roughness field redistributes storage rather than
throughput. The sign of the peak-depth difference also flips --- the
variable field gives the marginally lower maximum at baseline (4.36 vs
4.39 m) and the marginally higher one at the peak (4.57 vs 4.55 m) ---
so we draw no conclusion from a difference of two centimetres in a
single cell.

The reach is laterally confined, and it is worth stating what that means
for a two-dimensional solver rather than leaving a reader to infer it.
The wetted sheet is two cells wide at the median, which invites the
reading that a one-dimensional routing would serve equally well. The
measurements say otherwise about the cause and agree about the
consequence. The confinement is topographic, not a matter of the reach
being under-forced: across the rows that carry water, the valley floor
--- cells within 2 m of the thalweg --- is itself two cells wide at the
median, and the water occupies essentially all of it. Where the valley
opens to 300 m the inundation opens to 300 m with it. Nor would a larger
event change this: between the baseline and the peak the discharge rises
by a factor of 2.2 and the stored volume by 46 \%, while the wetted area
rises 6 \%, an empirical scaling of area with \(Q^{0.07}\) under which
even the 107 m$^{3}$/s maximum of the 92-year record --- 2.8 times the 2017
peak --- would not materially widen the flow.

What the solver is therefore doing here is tracking a valley floor whose
width varies from 60 to 300 m along the reach, which is information a
one-dimensional model would need supplied as surveyed cross-sections.
That is a genuine use of the second dimension, but a modest one, and at
30 m resolution a 60 m valley floor is two cells, which is marginal for
a finite-volume scheme. The benchmarks of §3 carry the harder
two-dimensional cases --- Test 8A is an urban configuration with
rainfall and surcharge at 2 m resolution, and the radial dam-break tests
isotropy directly. This application demonstrates the GIS-native input
pipeline and the roughness question, not two-dimensional hydrodynamics
at their most demanding.

This application is a demonstration of capability, not a calibrated
hindcast: the absolute depths depend on the warm-start and the
literature land-cover roughness, and a validated reconstruction would
require an official rating curve and field roughness survey. The point
is that the solver ingests standard public GIS products (SRTM DEM, ESA
WorldCover) and produces a physically coherent, mass-conserving
inundation field with spatially resolved friction --- the input pipeline
that a continental-scale, multi-basin programme requires.

\subsection{Sensitivity to the roughness
lookup}

The storage result rests on a literature lookup, so we sweep the three
land-cover classes that occupy the channel and its banks one at a time
over the range each carries in the standard compilations, holding the
other two at the §4.2 baseline and comparing every configuration against
the same uniform \(n = 0.04\) reference. The sweep is run at the day-one
baseline discharge of 17.5 m$^{3}$/s, where the reference comparison gives
+22 \% rather than the +29 \% of the peak; it is the \emph{robustness of
the direction} that transfers between discharges, not the magnitude, and
§4.3 has already shown the magnitude to be discharge-dependent.

\textbf{Table 3. One-at-a-time sensitivity of the roughness lookup, at
the day-one baseline discharge.} Baseline lookup is tree 0.100, shrub
0.060, bare 0.025; the reference is the uniform \(n = 0.04\) field.

\begin{longtable}[]{@{}
  >{\raggedright\arraybackslash}p{(\columnwidth - 8\tabcolsep) * \real{0.2000}}
  >{\raggedright\arraybackslash}p{(\columnwidth - 8\tabcolsep) * \real{0.2000}}
  >{\raggedright\arraybackslash}p{(\columnwidth - 8\tabcolsep) * \real{0.2000}}
  >{\raggedright\arraybackslash}p{(\columnwidth - 8\tabcolsep) * \real{0.2000}}
  >{\raggedright\arraybackslash}p{(\columnwidth - 8\tabcolsep) * \real{0.2000}}@{}}
\toprule\noalign{}
\begin{minipage}[b]{\linewidth}\raggedright
Swept class
\end{minipage} & \begin{minipage}[b]{\linewidth}\raggedright
\(n\)
\end{minipage} & \begin{minipage}[b]{\linewidth}\raggedright
Retained volume
\end{minipage} & \begin{minipage}[b]{\linewidth}\raggedright
Mean outflow
\end{minipage} & \begin{minipage}[b]{\linewidth}\raggedright
\(\overline{\Delta h}\) channel
\end{minipage} \\
\midrule\noalign{}
\endhead
\bottomrule\noalign{}
\endlastfoot
\emph{(baseline)} & --- & +22.3 \% & −3.6 \% & +0.187 m \\
Tree & 0.060 & +9.5 \% & −1.6 \% & +0.078 m \\
Tree & 0.150 & +38.2 \% & −6.3 \% & +0.322 m \\
Shrub & 0.040 & +21.6 \% & −3.5 \% & +0.181 m \\
Shrub & 0.080 & +23.4 \% & −3.8 \% & +0.199 m \\
Bare & 0.020 & +22.3 \% & −3.6 \% & +0.188 m \\
Bare & 0.030 & +22.3 \% & −3.7 \% & +0.189 m \\
\emph{all-minimum corner} & --- & +8.7 \% & −1.4 \% & +0.067 m \\
\emph{all-maximum corner} & --- & +39.1 \% & −6.4 \% & +0.331 m \\
\end{longtable}

A one-at-a-time design varies one class at a time and therefore cannot,
on its own, support a statement about the whole parameter box, so the
last two rows add the corners: all three classes at their minimum, and
all three at their maximum. The all-minimum corner is the adversarial
case, because it drives the channel classes toward the uniform
\(n = 0.04\) the comparison is against.

Three things follow. First, the \emph{direction} is robust: every
configuration --- including both corners --- retains more water, reduces
outflow, and deepens the channel relative to the uniform field. The sign
does not flip anywhere we sampled, corners included, so the qualitative
conclusion of §4.3 does not depend on the lookup. Second, the
\emph{magnitude} is not: the headline storage change spans +8.7 \% to
+39.1 \%, a more than fourfold range, and ``+22 \%'' should be read as
the midpoint of that band rather than as a determined quantity. Third,
and most usefully, essentially all of that spread is carried by a single
parameter. Varying the tree class moves the result across the whole
range; the shrub class moves it by under two percentage points; and the
bare class --- 66 \% of the domain by area, but hillslope rather than
channel --- moves it not at all, because the flow stays confined to the
channel at this discharge. The three classes were not swept over ranges
of equal relative width, so we also normalise: expressed as an
elasticity (fractional change in retained volume per fractional change
in \(n\)), the tree class returns \(\approx 0.26\) and the shrub class
\(\approx 0.02\), while the bare class falls below what this sweep
resolves --- both bare configurations return the same retained volume to
the precision of Table 3. These are two-point differences across the
swept ranges, so the figures carry only the precision that repetition
justifies: the tree value is quoted to two digits because the two
independent directions agree to two digits (0.26 decreasing, 0.26
increasing), which also indicates the response over this range is closer
to a power law in \(n\) than to a local linearisation, whereas the shrub
value is quoted to one. What the sweep establishes is the \emph{ranking}
--- an order of magnitude from tree to shrub and a further order to bare
--- and that ordering is not an artefact of the ranges having been
chosen with unequal relative width. The corners also bound the
interaction the one-at-a-time rows cannot see: moving shrub and bare to
their minima on top of tree's minimum shifts the result by only 0.8
percentage points (+9.5 \% to +8.7 \%), so the classes act very nearly
independently over this range.

That concentration matters for calibration strategy: the quantity that
needs constraining here is not a domain-averaged roughness but one
class-specific value, which places the problem on the favourable side of
the \(P^{*} \approx 2\) break-even of §2.5. Whether the forward-mode
gradient can actually be used to recover it depends on the integration
horizon rather than on the parameter count --- over a full simulated day
the tangent of this configuration is unstable (§2.5), so a gradient
calibration must work within a short assimilation window or against a
near-steady target. The companion study \citep{ParraPaper02} pursues
that line in one dimension, where the horizons are short enough for the
tangent to stay bounded. Reproduced by the
\texttt{huasco\_manning\_sweep} example.

\subsection{Cross-validation against an independent
solver}

The reach admits no observational validation of the inundation field:
there is no gauge inside the domain, the next station downstream sits
$\sim$35 km away in a valley where irrigation abstraction
dominates the difference between the two records, and a Sentinel-2 MNDWI
check across the event returned no detectable water signal --- the
valley carries about 5 m of relief across 360 m of width, below what a
10 m optical sensor separates. Rather than leave §4 with no external
check at all, we compare against an independent solver on the same
terrain.

\emph{SynxFlow} \citep{Xia2025} is a GPU-accelerated multi-hazard
package descended from HiPIMS-CUDA. It shares this solver's problem
class --- a well-balanced finite-volume discretisation of the 2D
shallow-water equations on a structured Cartesian grid --- while
differing in scheme, implementation language and hardware, which is what
makes the comparison informative. Both read GeoTIFF, so the same DEM
enters both codes unmodified. To remove the most obvious source of
drift, we do not let each code derive its own roughness: the Manning
field that hydroflux \emph{resolves} from the land cover is exported
cell-for-cell and supplied to SynxFlow as a gridded parameter, so the
two integrate identical friction.

One methodological point governs how the comparison must be set up, and
we report it because it would otherwise contaminate the result silently.
The two codes realise an upstream discharge differently: hydroflux
injects an exact volumetric point source, whereas SynxFlow converts a
discharge boundary series to velocities, so the delivered flux is the
product of velocity, depth and width and matches the target only insofar
as the depth used in the conversion is consistent. Driving both at 17.5
m$^{3}$/s, SynxFlow delivered 18.98 m$^{3}$/s effective --- 8.4 \% more water than
hydroflux received. Neither behaviour is wrong; they are two legitimate
ways to impose the same physical condition. But a comparison in which
one solver is given 8 \% more water measures the boundary treatment, not
the scheme.

We therefore seal both domains and remove the source, so that a full day
of integration redistributes an identical initial water body and any
difference is attributable to the numerics alone. On the 200 × 67-cell
reach at 30 m, over 86 400 s:

\textbf{Table 4. hydroflux versus SynxFlow, sealed domain, no source.}

\begin{longtable}[]{@{}ll@{}}
\toprule\noalign{}
Metric & Value \\
\midrule\noalign{}
\endhead
\bottomrule\noalign{}
\endlastfoot
RMSE in depth & 0.0210 m \\
MAE & 0.0097 m \\
Bias (SynxFlow − hydroflux) & +0.0002 m \\
Peak depth & 3.091 m vs 3.071 m \\
Critical Success Index, wet mask & 0.950 \\
Total stored volume & −0.022 \% \\
\end{longtable}

Two independent solvers agree to 2 cm RMSE, with a bias of 0.2 mm on
depths of order 3 m, on real 30 m terrain.

A single pair of runs cannot say whether that 2 cm is a genuine
scheme-level difference or an artefact of the time step each code
happened to use, so we vary ours. Repeating the comparison at CFL 0.4,
0.3 and 0.2 --- 78 000, 104 000 and 156 000 steps for the same simulated
day --- moves the RMSE from 0.0210 to 0.0208 to 0.0206 m, a spread of
1.9 \% across a doubling of the step count. The bias moves from +0.0002
to +0.0003 m and the peak depth from 3.071 to 3.070 m; the wet-mask CSI
goes from 0.950 to 0.945, which is one cell changing classification out
of the 199 the mask contains and is therefore the metric's granularity
rather than a trend. Expressed as a scaling, the residual falls with the
CFL number at an apparent order of 0.03: it does not scale with
\(\Delta t\).

That is the useful form of the answer. Had the disagreement been our
temporal discretisation, the apparent order would sit between one and
two and the residual would fall visibly under refinement; instead it
extrapolates to roughly 0.020 m rather than to zero as
\(\Delta t \to 0\). The 2 cm is therefore a scheme-level difference
between the two codes --- which is what the comparison is meant to
measure --- and not a consequence of having run at a loose CFL number.
Mass closure is unaffected across the three: \(4\cdot 10^{-15}\),
\(1.2\cdot 10^{-14}\) and \(6\cdot 10^{-15}\) relative. We vary only our
own time step here; the other code's discretisation settings are held at
their defaults, so this bounds our contribution to the residual rather
than decomposing it fully. For reference, the same comparison run with
the mismatched inflow gives an RMSE of 0.46 m and a bias of +0.31 m:
99.9 \% of the apparent disagreement was the injection mechanism rather
than the discretisation.

Mass behaviour is worth stating separately, because the sealed
configuration makes it exact rather than approximate. With the domain
closed and the point source active, hydroflux ends the day holding
\(1.564659\cdot 10^{6}\) m$^{3}$ against an analytically determined
\(1.564659\cdot 10^{6}\) m$^{3}$ --- a closure error of
\(-9.8\cdot 10^{-15}\) relative. This is a stricter test than the
closed-domain Thacker oscillation of §3.2, which conserves mass but has
no source term to balance. SynxFlow, integrated over the same sealed
domain with no source, drifts by \(-0.002\) \%.

What this establishes and what it does not both need saying. It
establishes that the redistribution machinery --- face fluxes, MUSCL
reconstruction, the wet/dry treatment and the friction step --- behaves
essentially identically to a mature, independently developed community
solver on real terrain. It does not validate the treatment of
\emph{sources} or of \emph{open} boundaries, which is precisely where
the two codes differ, and it is a model-to-model comparison rather than
an observational one: agreement between two solvers of the same problem
class is evidence against implementation error, not against shared
modelling assumptions. The clean configuration also runs without a
source, so it does not exercise the forced transient of §4.3. Full
numerical detail, including the isolation sequence that separated the
boundary artefact from the scheme, is in the repository.

\section{Roadmap}

The verified 2D solver is the foundation of a staged programme. The
immediate next layers are: (i) \textbf{GPU acceleration} via \(wgpu\)
compute shaders, prioritised because even the corrected, row-chunked CPU
parallelism (§3.9) saturates at \(3\)--\(4×\) on commodity multi-core
hardware --- well short of the throughput SIMT execution is expected to
provide on the same embarrassingly parallel face/cell loops; both the
well-balanced flux assembly and the cell-mask skip map naturally to
per-face/per-cell kernels, and the \texttt{Real}-trait dispatch shown to
work at no run-time cost in §2.5 is one of the structural assumptions
that should carry through unchanged on the GPU side. (ii) \emph{Physical
coupling} --- rainfall → slope-failure → granular propagation →
inundation in a single engine, beginning with an Iverson-type
debris-flow source \citep{Iverson2000, Christen2010} feeding the SW
momentum. (iii) \emph{Reverse-mode automatic differentiation}, required
to calibrate spatially distributed fields (per-cell roughness,
bathymetry corrections) whose parameter count far exceeds the measured
forward-mode break-even of \(P^{*} \approx 2\) (§2.5). The 1D companion
line already demonstrates forward-mode calibration of Manning and
cross-section parameters against real gauge data \citep{ParraPaper02};
the n--shape and n--bathymetry confounds it identifies motivate the
spatially distributed observations (remote-sensing inundation extent,
distributed stage) that the 2D solver is designed to assimilate. The
terminal goal is continental-scale coupled-hazard simulation across the
15 main Chilean basins, reproducible from versioned project files and
public data.

\section{Conclusion}

hydroflux is a 2D shallow-water finite-volume solver, written in Rust
and generic over its numeric type, that is well-balanced to machine
precision, mass-conservative on moving wet/dry fronts, and verified
against a hierarchy of analytical (lake-at-rest, Thacker, Stoker,
MacDonald) and community (UK EA: 6 synthetic stand-ins, plus Tests 4 and
8A on the official EA/LISFLOOD-FP geometry) benchmarks. It ingests
standard public GIS products directly and, applied to a Río Huasco reach
forced by a regulated 2017 reservoir release, demonstrates a
land-cover-derived spatially variable Manning field that, as a
sensitivity test rather than a validated hindcast, retains
$\sim$29 \% more water in the reach at the peak of the routed
release than a single uniform roughness --- a direction that survives
every configuration of a roughness sweep, corners of the parameter box
included. On the same terrain the solver reproduces the depth field of
an independently developed GPU community solver to 0.021 m RMSE. The
contribution is a verified, open artifact --- the
differentiable-by-design 2D foundation of a multi-year programme toward
coupled, GPU-native, continental-scale hazard simulation --- released
for the community to build on.



%% ============================================================
%% Figures — embedded at end of body, before bibliography.
%% Captions mirror the Markdown source; the \ref labels match the
%% sed-applied transforms in the Makefile that turn "Figure N" into
%% \ref{fig:...} in the body text.
%% ============================================================

\begin{figure}[!htbp]
\centering
\includegraphics[width=0.85\textwidth]{figures/fig01_scheme.pdf}
\caption{Well-balanced finite-volume scheme at an x-face between cells L and R over a stepped bed (to scale). The cells carry actual depths $h_L, h_R$ (blue brackets at the cell centres) with free surfaces $\eta _L, \eta _R$. The Audusse (2004) hydrostatic reconstruction defines $z_max = max(z_L, z_R)$ (purple dashed) and the reconstructed face depths $h*_L = max(\eta _L - z_max, 0)$, $h*_R = max(\eta _R - z_max, 0)$ (purple brackets), on which the HLLC flux $F$ (vermillion) is evaluated. The bed is reconstructed linearly to the shared face value $z_face = ½(z_L + z_R)$, which makes the Audusse pressure correction vanish and moves the bed-slope force into the cell-centred source term (orange) --- the construction that gives the exact lake-at-rest balance verified in \S3.1. Schematic (no simulation data); generated by \texttt{fig01\_scheme.R}.}
\label{fig:scheme}
\end{figure}

\begin{figure}[!htbp]
\centering
\includegraphics[width=\textwidth]{figures/fig02_verification.pdf}
\caption{Verification against analytical references: simulated profiles (points) overlaid on the analytical solution (line). (a) Stoker/Ritter dam-break on a dry bed at $t = 4$ s --- the rarefaction fan and the leading dry front are captured (400-cell channel, mid-row slice). (b) MacDonald uniform flow at steady state, y- axis zoomed to ±1.5 \% of the normal depth \texttt{h\_n} to expose the well-balanced preservation (the profile stays flat at \texttt{h\_n} on a sloped bed under Manning friction). (c) Thacker oscillating paraboloid at $t = T/2$, centre-row slice --- the parabolic cap and its moving wet/dry shoreline track the analytical solution on a curved bed. Quantitative error metrics for all benchmarks are in Table 1 (computed over the benchmark-specific regions: the Stoker rarefaction fan, the 2D Thacker interior); they are not annotated here to avoid conflating the 1D-slice view of this figure with the table's domain-wide norms. Generated by \texttt{fig02\_verification.R} from the \texttt{gen\_verification\_data} example output.}
\label{fig:verification}
\end{figure}

\begin{figure}[!htbp]
\centering
\includegraphics[width=\textwidth]{figures/fig03_uk_ea_t6.pdf}
\caption{UK Environment Agency 2D benchmark Test 6: urban dam-break depth field at $t = 30$ s on the 500 $\times$ 100 m domain (250 $\times$ 50 cells, $\Delta x = \Delta y = 2$ m). The reservoir ($x < 100$ m, dashed dam line) collapses into the dry downstream and reaches the first building cluster ($x \approx 150$--170 m) by $t = 30$ s, wrapping around the obstacles with the expected wakes and inter-building jets. The six raised-bed buildings (grey) remain dry throughout --- water navigates the array without overtopping them --- while the wet/dry front is captured without spurious oscillation. Depth on a perceptually uniform sequential colour scale (scico *devon*); downstream buildings ($x \geq 250$ m) are still dry at this time as the front has not yet arrived. Generated by \texttt{fig03\_uk\_ea\_t6.R} from the \texttt{gen\_verification\_data} example output.}
\label{fig:ukea_t6}
\end{figure}

\begin{figure}[!htbp]
\centering
\includegraphics[width=\textwidth]{figures/fig04_huasco_application.pdf}
\caption{Río Huasco reach under the regulated 2017 release, 200 $\times$ 67-cell 30 m subset (UTM 19S), at the day-11 peak. Five panels sharing the reach extent: (a) ESA WorldCover 2021 land cover --- riparian tree and shrub vegetation tracks the thalweg within bare Atacama hillslopes; (b) the derived Manning field $n(x, y)$ mapping land cover to roughness ($n = 0.025$ bare to $0.10$ tree); (c) inundation depth with a single uniform $n = 0.04$; (d) inundation depth with the variable $n(x, y)$; (e) the difference $\Delta h = (d) - (c)$, hill-shaded base, divergent scale. Panels (c)--(e) share a hillshade underlay. The positive $\Delta h$ (warm) concentrated in the channel shows the riparian roughness deepening and retaining the flow --- the +0.33 m mean channel deepening and +29 \% retained volume reported in \S4.3. Generated by \texttt{fig04\_huasco\_application.R}, which reuses the solver-2d example rasters (\texttt{huasco\_subset\_{dem,landcover}.tif}, \texttt{huasco\_2d\_depth\_day\_11{,\_landcover}.tif}).}
\label{fig:huasco}
\end{figure}

\begin{figure}[!htbp]
\centering
\includegraphics[width=0.75\textwidth]{figures/fig05_convergence.pdf}
\caption{Mesh-refinement convergence on the Thacker oscillation (\S3.7). (a) Relative L1 (circles) and L2 (triangles) error in depth at $t = T/2$ versus cell size $\Delta x$, log-log, with order-1 and order-2 reference slopes (dashed); the overall fitted slopes are 1.73 (L1) and 1.58 (L2). (b) Order observed between successive meshes, $\log_2(e_{\text{coarse}}/e_{\text{fine}})$, plotted against the finer mesh of each pair. Panel (b) carries the claim panel (a) cannot: a single fitted slope averages the degradation away, whereas the pairwise orders show it directly, falling from 1.81 to 1.59 in L1 and from 1.71 to 1.43 in L2 as the moving wet/dry shoreline --- locally first order --- comes to dominate the shrinking smooth-region error. Generated by \texttt{fig05\_convergence.R} from the \texttt{gen\_convergence} example.}
\label{fig:convergence}
\end{figure}

\begin{figure}[!htbp]
\centering
\includegraphics[width=0.9\textwidth]{figures/fig06_head_to_head.pdf}
\caption{Head-to-head against ANUGA \citep{Roberts2015} on the Stoker/Ritter dam-break at matched resolution ($\Delta x = 1$ m, \S3.8). Both solvers (circles: hydroflux; triangles: ANUGA) overlaid on the analytical Ritter profile (line) at $t = 4$ s; the rarefaction fan $[37.5, 75.1]$ m is shaded. The dashed vertical marks the initial dam at $x = 50$ m. Relative error norms over the fan, annotated: L1 4.1 \% / L2 3.6 \% / $L^{\infty}$ 5.3 \% $h_L$ for hydroflux versus 2.6 \% / 2.7 \% / 4.4 \% for ANUGA. Both stay below 6 \% at this coarse resolution; hydroflux's reference 400-cell run (Table 1) reaches L1 1.0 \%. Generated by \texttt{fig06\_head\_to\_head.R} from the \texttt{anuga\_stoker\_compare.py} + \texttt{gen\_stoker\_coarse} example outputs.}
\label{fig:head_to_head}
\end{figure}


\section*{Author Contributions}

CRediT roles \citep{Brand2015}:

\begin{itemize}
\item \textbf{Francisco Parra} --- Conceptualization; Methodology;
  Software (solver-2d, autograd, verification suite, application
  pipeline); Validation; Formal analysis; Investigation; Data
  curation; Writing --- original draft; Writing --- review \&
  editing; Visualization; Project administration; Funding acquisition
  (DICYT postdoctoral fellowship).
\item \textbf{Verónica Gil-Costa} --- Methodology (numerical schemes
  and benchmarking); Validation; Writing --- review \& editing;
  Supervision (computational methodology).
\item \textbf{Carolina Bonacic} --- Software (high-performance
  computing and parallelism review); Validation; Writing --- review
  \& editing.
\item \textbf{Mauricio Marín} --- Conceptualization (research line
  and venue positioning); Methodology (software-engineering rigor
  and reproducibility); Writing --- review \& editing; Supervision
  (Rust systems programming and software engineering); Resources
  (computing infrastructure).
\end{itemize}

All authors read and approved the final manuscript.

\section*{Open Research}

The solver source code (Rust), the test suite, all benchmark
configurations (lake-at-rest, Thacker, Stoker, MacDonald, radial
dam-break, UK EA tests 1--6, the mesh-refinement convergence study,
the head-to-head versus ANUGA, and the M1 / M2 examples for §2.5 and
§3.9), the Huasco application notebook, and the figure-generation
scripts are released under the MIT OR Apache-2.0 dual licence at
\url{https://github.com/franciscoparrao/hydroflux}. A Zenodo snapshot
DOI of the version associated with this manuscript will be minted at
acceptance and added to this section in proof.

The ESA WorldCover 2021 land-cover raster used in §4 is available at
\url{https://esa-worldcover.org/} \citep{ESAWorldCover2021}. The
30~m SRTM DEM is publicly available from the USGS. The DGA gauge
record at Santa Juana (03820003) is available on request from the
Dirección General de Aguas, Chile.

\section*{Declaration of competing interest}

The author declares no competing financial or personal interests.

\section*{Acknowledgements}

This work was supported by the Postdoctoral DICYT 2026--2027 programme
of the Universidad de Santiago de Chile.

\section*{Declaration on generative AI use}

Generative AI tools were used to assist drafting, refactoring, and
reference verification. All content was reviewed and validated by the
author. Three references in an earlier draft (Hydrograd.jl, AegirJAX,
r.avaflow v4) failed verify-refs / OpenAlex checks and were removed
before submission.

\bibliographystyle{elsarticle-harv}
\bibliography{references}

\end{document}
